\begin{center}
\LARGE
The Role of Representation for Specification and
          Communication in Mathematical Assistants

\end{center}

\begin{center}
\Large
Jacques Calmet, Karsten Homann

\end{center}

\noindent
Date:  July 17th (Wednesday)\\
Time:  16:05-16:30\\

\begin{center}
\large
Abstract
\end{center}

It is well known that mathematicians switch to different views of a problem when needed. They
can represent a problem at a formal, conceptual, heuristic, algorithmic or constraint level
whenever necessary. To represent Mathematics within several formalisms has been the subject of
many research projects. However, only few knowledge-based systems manage translations of
representations between theories. 

The goal of this talk is twofold. One the one hand, we report on a hybrid knowledge representation
and reasoning system called MANTRA. The system provides four different knowledge
representation methods - first-order logic, terminological language, semantic networks, and
production rules - distributed into a three levels architecture. Specifications of mathematical
domains of computation and their inherently related type inference mechanisms can be
transformed into knowledge bases. 

On the other hand, we argue that a main requirement when designing future environments is the
capability to cooperate and to integrate by communicating mathematical knowledge
among/through mathematical services based upon restart-able computation and reasoning.
Therefore, structures representing intermediate results in any kind of mathematical computation
must also be considered. 

